\documentclass[12pt]{article}

\usepackage[paperwidth=612pt,paperheight=792pt,top=72pt,bottom=72pt,left=30mm,right=20mm]{geometry}
\usepackage{times}
\usepackage{setspace}
\usepackage{fancyhdr}
\usepackage{ragged2e}
\usepackage{graphicx}
\usepackage{mdframed}
\usepackage{listings}
\usepackage{color}
\setstretch{1.5}

\definecolor{codegreen}{rgb}{0,0.6,0}
\definecolor{codegray}{rgb}{0.5,0.5,0.5}
\definecolor{codepurple}{rgb}{0.58,0,0.82}
\definecolor{backcolour}{rgb}{0.95,0.95,0.92}

\lstdefinestyle{mystyle}{
    backgroundcolor=\color{backcolour},
    commentstyle=\color{codegreen},
    keywordstyle=\color{magenta},
    numberstyle=\tiny\color{codegray},
    stringstyle=\color{codepurple},
    basicstyle=\ttfamily\footnotesize,
    breakatwhitespace=false,
    breaklines=true,
    captionpos=b,
    keepspaces=true,
    numbers=left,
    numbersep=5pt,
    showspaces=false,
    showstringspaces=false,
    showtabs=false,
    tabsize=2
}
\lstset{style=mystyle}

\pagestyle{fancy}
\fancyhf{}
\fancyhead[R]{AyurSpace}
\fancyfoot[L]{Pillai HOC College of Engineering \& Technology (Autonomous), B. E. Computer Engineering}
\fancyfoot[R]{\thepage}
\renewcommand{\headrulewidth}{1pt}
\renewcommand{\footrulewidth}{1pt}

\begin{document}

\vspace*{\fill}
\begin{center}
\textbf{\Huge{Chapter 7}}\\
\vspace{5mm}
\textbf{\Huge{Implemented System}}
\end{center}
\vspace*{\fill}
\newpage

\section*{7.1 \quad System Architecture}
\addcontentsline{toc}{section}{7.1 System Architecture}

\renewcommand{\thefigure}{7.1}
\renewcommand{\figurename}{Figure}
\begin{figure}[!htbp]
\begin{mdframed}
    \centering
    \includegraphics[width=0.7\linewidth,keepaspectratio]{SystemArch.png}
\end{mdframed}
\caption{Implemented System Architecture}
\end{figure}

\justify
The central infrastructure revolves around solving complex workflows via independent processes.
High-resolution images from standard modern smartphone sensors (12MP) incur payload cost in network
transmissions. A local Pre-processing function $P(I_{raw})$ performs bicubic downsampling, capping
dimensions and heavily lossy compressing to JPEG Q=85. This decreases payloads by 84\% before
interfacing with REST APIs.

Our dual-stage API pipeline (Hybrid Inference) is built such that Plant.id supplies visual
confidence, effectively serving as the discriminative boundary for the Google Gemini payload. By
employing specific System Prompts acting as `AyurBot', we curb any ``Alignment Drift'' that general
LLMs hold regarding western interpretation of medicinal texts. The state updates in the Flutter UI
smoothly execute over independent asynchronous threads connected via Riverpod. This methodology
ensures AyurSpace retains efficiency on affordable hardware seen throughout various demographics.

\newpage

\section*{7.2 \quad State Management with Riverpod}
\addcontentsline{toc}{section}{7.2 State Management with Riverpod}
\justify

AyurSpace heavily utilizes Riverpod for reactive state management and dependency injection.
Riverpod provides a compile-time safe mechanism to access global state without the pitfalls of
standard InheritedWidgets or Provider. The Chat functionality requires maintaining a history of
messages, handling loading states (typing indicators), and communicating with the Gemini service.

The \texttt{ChatState} encompasses the list of messages, the typing status, and any potential
errors encountered during the AI query process:

\begin{lstlisting}[language=C++]
class ChatState {
  final List<ChatMessage> messages;
  final bool isTyping;
  final String? error;

  const ChatState({
    this.messages = const [],
    this.isTyping = false,
    this.error,
  });

  ChatState copyWith({
    List<ChatMessage>? messages,
    bool? isTyping,
    String? error,
    bool clearError = false,
  }) {
    return ChatState(
      messages: messages ?? this.messages,
      isTyping: isTyping ?? this.isTyping,
      error: clearError ? null : (error ?? this.error),
    );
  }
}
\end{lstlisting}

The system contains multiple providers corresponding to specific domains of the application:
\begin{itemize}
    \item \textbf{AuthProvider}: Manages Firebase Authentication states (logged in, logged out,
    loading).
    \item \textbf{ScanProvider}: Manages the image capture, upload, and processing pipeline for
    plant identification.
    \item \textbf{DoshaQuizProvider}: Manages the step-by-step questionnaire state to calculate
    the user's \textit{Prakriti}.
    \item \textbf{UserProvider}: Syncs the user's profile and preferences with Cloud Firestore.
\end{itemize}

\newpage

\section*{7.3 \quad External AI Integrations}
\addcontentsline{toc}{section}{7.3 External AI Integrations}

\subsection*{7.3.1 \quad API Configuration}
The application centralizes its API keys and endpoints in a secure configuration module, parsed
from a \texttt{.env} file.

\begin{lstlisting}[language=C++]
import 'package:flutter_dotenv/flutter_dotenv.dart';

class ApiConfig {
  ApiConfig._();

  static String get plantIdApiKey =>
      dotenv.env['PLANT_ID_API_KEY'] ?? '';
  static const String plantIdBaseUrl = 'https://plant.id/api/v3';

  static String get geminiApiKey =>
      dotenv.env['GEMINI_API_KEY'] ?? '';
  static const String geminiModel = 'gemini-2.5-flash';

  static bool get isPlantIdConfigured =>
      plantIdApiKey.isNotEmpty &&
      plantIdApiKey != 'YOUR_PLANT_ID_API_KEY';
  static bool get isGeminiConfigured =>
      geminiApiKey.isNotEmpty &&
      geminiApiKey != 'YOUR_GEMINI_API_KEY';
}
\end{lstlisting}

\subsection*{7.3.2 \quad Plant.id Integration}
The \texttt{PlantIdService} communicates with the Kindwise Plant.id API using REST over HTTP. The
request requires the base64 encoded image sequence and configuration options dictating what data to
return (common names, descriptions, edibility, propagation methods). The JSON response from the API
is deeply nested. The service maps this JSON to the \texttt{PlantIdResult} Dart object, abstracting
away the JSON structure from the rest of the application.

\subsection*{7.3.3 \quad Google Gemini Integration}
AyurSpace heavily relies on the \texttt{gemini-2.5-flash} model due to its high inference speed
and robust multimodal capabilities. The \texttt{GeminiService} has distinct methods depending on
the operation:
\begin{itemize}
    \item \texttt{getAyurvedicInfo()}: Zero-shot prompt for retrieving data about a specific
    identified plant.
    \item \texttt{sendChat()}: Multi-turn conversation utilizing the \texttt{system\_instruction}
    payload configuration alongside the historical \texttt{contents} array.
    \item \texttt{analyzeImage()}: Multimodal prompt combining \texttt{inline\_data} image payload
    and a text prompt to perform end-to-end analysis.
\end{itemize}

Safety settings are injected into every Gemini request to prevent generation of harmful content,
specifically filtering for Hate Speech, Dangerous Content, Harassment, and Sexually Explicit
material.

\newpage

\section*{7.4 \quad Complete Source Code Listings}
\addcontentsline{toc}{section}{7.4 Complete Source Code Listings}
\justify

\textbf{Plant.id Service Implementation (\texttt{plant\_id\_service.dart}):}
\begin{lstlisting}[language=C++]
import 'dart:convert';
import 'package:flutter/foundation.dart';
import 'package:http/http.dart' as http;
import '../config/api_config.dart';

class PlantIdResult {
  final String scientificName;
  final String commonName;
  final double probability;
  final String? description;
  final String? imageUrl;
  final List<String> similarImages;
  final bool isHealthy;
  final String? healthAssessment;

  const PlantIdResult({
    required this.scientificName,
    required this.commonName,
    required this.probability,
    this.description,
    this.imageUrl,
    this.similarImages = const [],
    this.isHealthy = true,
    this.healthAssessment,
  });

  factory PlantIdResult.fromJson(Map<String, dynamic> json) {
    final suggestion =
        json['result']['classification']['suggestions'][0];
    final details = suggestion['details'] ?? {};

    return PlantIdResult(
      scientificName: suggestion['name'] ?? 'Unknown',
      commonName: details['common_names']?.first
          ?? suggestion['name'] ?? 'Unknown',
      probability:
          (suggestion['probability'] ?? 0.0).toDouble(),
      description: details['description']?['value'],
      imageUrl: details['image']?['value'],
      similarImages: List<String>.from(
        (suggestion['similar_images'] ?? [])
            .map((img) => img['url'] ?? ''),
      ),
      isHealthy:
          json['result']['is_healthy']?['binary'] ?? true,
      healthAssessment:
          json['result']['is_healthy']?['assessment'],
    );
  }
}

class PlantIdService {
  final http.Client _client;

  PlantIdService({http.Client? client})
      : _client = client ?? http.Client();

  Future<PlantIdResult> identifyFromBytes(
      Uint8List imageBytes) async {
    if (!ApiConfig.isPlantIdConfigured) {
      throw PlantIdException('API not configured');
    }
    try {
      final base64Image = base64Encode(imageBytes);
      final response = await _client.post(
        Uri.parse(
            '${ApiConfig.plantIdBaseUrl}/identification'),
        headers: {
          'Api-Key': ApiConfig.plantIdApiKey,
          'Content-Type': 'application/json',
        },
        body: jsonEncode({
          'images': [
            'data:image/jpeg;base64,$base64Image'
          ],
          'similar_images': true,
          'health': 'auto',
          'language': 'en',
          'details': [
            'common_names', 'description',
            'image', 'edible_parts',
            'watering', 'propagation_methods',
          ],
        }),
      );
      if (response.statusCode == 200 ||
          response.statusCode == 201) {
        final json = jsonDecode(response.body);
        return PlantIdResult.fromJson(json);
      } else {
        throw PlantIdException(
            'API error: ${response.statusCode}');
      }
    } catch (e) {
      rethrow;
    }
  }

  void dispose() => _client.close();
}

class PlantIdException implements Exception {
  final String message;
  const PlantIdException(this.message);
  @override
  String toString() => 'PlantIdException: $message';
}
\end{lstlisting}

\vspace{1cm}

\textbf{Gemini Generative Flow Implementation (\texttt{gemini\_service.dart}):}
\begin{lstlisting}[language=C++]
import 'dart:convert';
import 'package:http/http.dart' as http;
import '../config/api_config.dart';

class GeminiResponse {
  final String text;
  final bool isError;
  const GeminiResponse({required this.text,
      this.isError = false});
}

class ChatTurn {
  final String role;
  final String text;
  const ChatTurn({required this.role, required this.text});
  Map<String, dynamic> toJson() => {
    'role': role,
    'parts': [{'text': text}],
  };
}

class GeminiService {
  final http.Client _client;
  GeminiService({http.Client? client})
      : _client = client ?? http.Client();

  static const List<Map<String, String>> _safetySettings = [
    {'category': 'HARM_CATEGORY_DANGEROUS_CONTENT',
     'threshold': 'BLOCK_ONLY_HIGH'},
    {'category': 'HARM_CATEGORY_HARASSMENT',
     'threshold': 'BLOCK_ONLY_HIGH'},
    {'category': 'HARM_CATEGORY_HATE_SPEECH',
     'threshold': 'BLOCK_ONLY_HIGH'},
    {'category': 'HARM_CATEGORY_SEXUALLY_EXPLICIT',
     'threshold': 'BLOCK_ONLY_HIGH'},
  ];

  static const Map<String, dynamic> _generationConfig = {
    'temperature': 0.8,
    'maxOutputTokens': 2048,
    'topP': 0.95,
  };

  Future<GeminiResponse> getAyurvedicInfo({
    required String plantName,
    String? scientificName,
  }) async {
    final prompt = '''
You are AyurBot, a friendly Ayurvedic wellness guide.
Tell me about: **$plantName**
${scientificName != null
    ? '(Scientific name: *$scientificName*)' : ''}

Include:
- Hindi Name
- Dosha Effect
- Ayurvedic Uses
- Health Benefits
- How to Use at Home
- Precautions

Keep it friendly and under 300 words.
''';
    return await sendMessage(prompt);
  }

  Future<GeminiResponse> sendMessage(String message) async {
    if (!ApiConfig.isGeminiConfigured) {
      return const GeminiResponse(
          text: 'Mock response: API key not configured.');
    }
    try {
      final url =
          'https://generativelanguage.googleapis.com'
          '/v1beta/models/${ApiConfig.geminiModel}'
          ':generateContent';
      final response = await _client.post(
        Uri.parse(url),
        headers: {
          'Content-Type': 'application/json',
          'x-goog-api-key': ApiConfig.geminiApiKey,
        },
        body: jsonEncode({
          'contents': [{'parts': [{'text': message}]}],
          'generationConfig': _generationConfig,
          'safetySettings': _safetySettings,
        }),
      );
      if (response.statusCode == 200) {
        final json = jsonDecode(response.body);
        final candidates = json['candidates'] as List?;
        if (candidates != null && candidates.isNotEmpty) {
          final text = candidates[0]['content']
              ['parts'][0]['text'] as String?;
          if (text != null) return GeminiResponse(text: text);
        }
        return const GeminiResponse(
            text: 'Empty response', isError: true);
      }
      return const GeminiResponse(
          text: 'API Error', isError: true);
    } catch (e) {
      return const GeminiResponse(
          text: 'Connection Error', isError: true);
    }
  }

  void dispose() => _client.close();
}
\end{lstlisting}

\end{document}
